\documentclass[a4paper,12pt]{article}
\usepackage[utf8]{inputenc}
\usepackage[T1]{fontenc}
\usepackage[ngerman]{babel}
\usepackage{amsmath, amssymb}
\usepackage{geometry}
\usepackage{tikz}
\usepackage{tikz-cd}
\usetikzlibrary{shapes.geometric, arrows.meta, positioning, fit, backgrounds}
\usepackage{parskip}
\usepackage{titlesec}
\usepackage{booktabs}
\usepackage{enumitem}

\geometry{left=3cm, right=3cm, top=2.5cm, bottom=2.5cm}

\titleformat{\section}{\large\bfseries}{}{0em}{}[\titlerule]
\titleformat{\subsection}{\normalsize\bfseries}{}{0em}{}

\tikzset{
  box/.style={rectangle, draw, rounded corners=3pt, text width=4.5cm,
              align=center, minimum height=1.0cm, font=\small, inner sep=4pt},
  boxw/.style={box, text width=5cm},
  arrow/.style={-{Latex[length=2mm]}, thick},
}

\begin{document}

\begin{center}
  {\LARGE\bfseries Studierhinweise zu Lektion 2}\\[0.5em]
  {\large Mathematische Grundlagen (Modul 61111)}
\end{center}

\vspace{1em}

In dieser Lektion beginnt die eigentliche Lineare Algebra. Das zentrale Handwerkzeug
ist der \textbf{Gaußalgorithmus} (Kapitel~4), mit dem Matrizen in Treppennormalform
überführt werden. Er bildet die Grundlage für das Lösen linearer Gleichungssysteme
(Kapitel~5) und für den abstrakten Begriff des Vektorraums (Kapitel~6), der das gesamte
weitere Studium durchzieht.

Besondere Aufmerksamkeit verdienen der Gaußalgorithmus und die Lösungsformel
für lineare Gleichungssysteme – beide müssen Sie \emph{aktiv} beherrschen, nicht
nur passiv verstehen. Rechnen Sie viele Beispiele, nutzen Sie die Maple-Apps in der
Moodle-Lernumgebung, und arbeiten Sie die Aufgaben im Lehrtext durch.

% -----------------------------------------------------------------------
\section*{Struktur der Lektion 2}

\begin{center}
\begin{tikzpicture}[node distance=0.6cm and 0.8cm]

  % Column headers
  \node[font=\bfseries\small] (h1) at (0,0)    {Kapitel 4: Gaußalgorithmus};
  \node[font=\bfseries\small] (h2) at (6.5,0)  {Kapitel 5: Lin. Gleichungssysteme};
  \node[font=\bfseries\small] (h3) at (12,0) {Kapitel 6: Vektorräume};

  % Chapter 4 boxes
  \node[box, below=0.8cm of h1] (tnf)  {4.1 Treppennormalform\\Definition, Pivot-Positionen};
  \node[box, below=0.5cm of tnf] (gau) {4.2 Gaußalgorithmus\\spaltenweise Umformung};
  \node[box, below=0.5cm of gau] (tra) {4.3 Transformationsmatrix\\$\bigl(A\mid I_m\bigr)\to\bigl(T\mid S\bigr)$};
  \node[box, below=0.5cm of tra] (ein) {4.4 Eindeutigkeit\\der Treppennormalform};
  \node[box, below=0.5cm of ein] (rng) {4.5 Rang $\operatorname{Rg}(A)$\\Invertierbarkeit};

  % Chapter 5 boxes
  \node[box, below=0.8cm of h2] (lgs) {5.1 Definitionen\\$Ax=b$, Koeffizientenmatrix};
  \node[box, below=0.5cm of lgs] (str) {5.2 Struktur der Lösungsmenge\\$L = \lambda_0 + U$};
  \node[box, below=0.5cm of str] (los) {5.3 Zusammenfassung\\Merkregeln für $\lambda_0$ und $U$};

  % Chapter 6 boxes
  \node[box, below=0.8cm of h3] (vrd) {6.1 Vektorraum\\Axiome, Beispiele};
  \node[box, below=0.5cm of vrd] (unt) {6.2 Unterräume\\Unterraumkriterium, $U+W$, $U\cap W$};
  \node[box, below=0.5cm of unt] (erz) {6.3 Endlich erzeugte Vektorräume\\Linearkombination, Erzeugnis};

  % Arrows within Kap 4
  \draw[arrow] (tnf) -- (gau);
  \draw[arrow] (gau) -- (tra);
  \draw[arrow] (tra) -- (ein);
  \draw[arrow] (ein) -- (rng);

  % Arrows within Kap 5
  \draw[arrow] (lgs) -- (str);
  \draw[arrow] (str) -- (los);

  % Arrows within Kap 6
  \draw[arrow] (vrd) -- (unt);
  \draw[arrow] (unt) -- (erz);

  % Cross-chapter arrows
  \draw[arrow] (rng.east) -- ++(0.3,0) |- (lgs.west);
  \draw[arrow] (str.east) -- ++(0.3,0) |- (vrd.west);

\end{tikzpicture}
\end{center}

\newpage

% -----------------------------------------------------------------------
\section*{Zielelement 4.1 -- Treppennormalformen}

\subsection*{Lerninhalte}

\begin{center}
\begin{tikzpicture}[node distance=0.6cm and 1.0cm]
  \node[box] (def) {Definition der Treppennormalform\\(vier Bedingungen)};
  \node[box, below=0.5cm of def] (piv) {Pivot-Positionen $(i,j_i)$};
  \node[box, below=0.5cm of piv] (bsp) {Beispiele und Gegenbeispiele\\in $M_{mn}(K)$};
  \node[box, below=0.5cm of bsp] (inv) {Invertierbare Matrizen\\in Treppennormalform $\Leftrightarrow$ $I_m$};
  \draw[arrow] (def) -- (piv);
  \draw[arrow] (piv) -- (bsp);
  \draw[arrow] (bsp) -- (inv);
\end{tikzpicture}
\end{center}

\subsection*{Lernziele}

Nach Durcharbeiten dieses Abschnitts sollten Sie:
\begin{itemize}
  \item die vier Bedingungen der Treppennormalform kennen und auf konkrete Matrizen
        anwenden können,
  \item entscheiden können, ob eine gegebene Matrix in Treppennormalform ist,
        und gegebenenfalls die Pivot-Positionen angeben,
  \item wissen, dass die einzige invertierbare Matrix in Treppennormalform die
        Einheitsmatrix ist.
\end{itemize}

\subsection*{Selbstkontrollelement 4.1}

Entscheiden Sie, welche der folgenden Matrizen in Treppennormalform sind, und
geben Sie jeweils die Pivot-Positionen an:
\[
  A = \begin{pmatrix} 1 & 2 & 0 \\ 0 & 1 & 3 \\ 0 & 0 & 0 \end{pmatrix}, \quad
  B = \begin{pmatrix} 0 & 1 & 0 \\ 0 & 0 & 1 \\ 0 & 0 & 1 \end{pmatrix}, \quad
  C = \begin{pmatrix} 1 & 0 & 0 \\ 0 & 0 & 1 \\ 0 & 1 & 0 \end{pmatrix}.
\]

\newpage

% -----------------------------------------------------------------------
\section*{Zielelement 4.2 -- Der Gaußalgorithmus}

\subsection*{Lerninhalte}

\begin{center}
\begin{tikzpicture}[node distance=0.6cm and 1.0cm]
  \node[box] (idee) {Idee: spaltenweises Vorgehen};
  \node[box, below=0.5cm of idee] (fal) {Drei Fälle pro Spalte:\\kein Pivot / Pivot direkt / Zeilentausch};
  \node[box, below=0.5cm of fal] (alg) {Formaler Algorithmus\\(Satz 4.2.4)};
  \node[box, below=0.5cm of alg] (tra) {Transformationsmatrix $S$\\via $\bigl(A\mid I_m\bigr)$};
  \draw[arrow] (idee) -- (fal);
  \draw[arrow] (fal) -- (alg);
  \draw[arrow] (alg) -- (tra);
\end{tikzpicture}
\end{center}

\subsection*{Lernziele}

Nach Durcharbeiten dieses Abschnitts sollten Sie:
\begin{itemize}
  \item den Gaußalgorithmus sicher und fehlerfrei auf beliebige Matrizen anwenden können,
  \item die Transformationsmatrix $S$ mit $SA = T$ (Treppennormalform) berechnen,
  \item wissen, dass die Treppennormalform eindeutig ist (Satz~4.4.2), das Produkt
        $S$ der Elementarmatrizen aber nicht.
\end{itemize}

\subsection*{Selbstkontrollelement 4.2}

Überführen Sie die Matrix
\[
  A = \begin{pmatrix} 0 & 2 & -4 \\ 1 & 3 & -1 \\ 2 & 4 & 0 \end{pmatrix}
\]
durch elementare Zeilenumformungen in Treppennormalform und berechnen Sie dabei
gleichzeitig die invertierbare Matrix $S$ mit $SA = T$.

\newpage

% -----------------------------------------------------------------------
\section*{Zielelement 4.5 -- Rang und Invertierbarkeit}

\subsection*{Lerninhalte}

\begin{center}
\begin{tikzpicture}[node distance=0.6cm and 1.0cm]
  \node[box] (rng) {Definition $\operatorname{Rg}(A)$:\\Anzahl der Pivot-Positionen};
  \node[box, below=0.5cm of rng] (inv) {Proposition 4.5.4:\\(a)--(d) äquivalent};
  \node[box, below=0.5cm of inv] (kol) {Korollar 4.5.6:\\$CA = I_m \Rightarrow C = A^{-1}$};
  \node[box, below=0.5cm of kol] (ber) {Berechnung von $A^{-1}$\\via Gaußalgorithmus};
  \draw[arrow] (rng) -- (inv);
  \draw[arrow] (inv) -- (kol);
  \draw[arrow] (kol) -- (ber);
\end{tikzpicture}
\end{center}

\subsection*{Lernziele}

Nach Durcharbeiten dieses Abschnitts sollten Sie:
\begin{itemize}
  \item den Rang einer Matrix berechnen können,
  \item die vier äquivalenten Charakterisierungen invertierbarer Matrizen
        (Proposition~4.5.4) kennen und einsetzen können,
  \item Korollar~4.5.6 verstehen: aus $CA = I_m$ folgt bereits die volle Invertierbarkeit,
  \item die Inverse einer invertierbaren Matrix mit dem Gaußalgorithmus berechnen.
\end{itemize}

\subsection*{Selbstkontrollelement 4.5}

Untersuchen Sie, ob $A = \begin{pmatrix} 2 & 1 \\ 4 & 3 \end{pmatrix}$ invertierbar ist,
und berechnen Sie gegebenenfalls $A^{-1}$ mit dem Gaußalgorithmus.

\newpage

% -----------------------------------------------------------------------
\section*{Zielelement 5 -- Lineare Gleichungssysteme}

\subsection*{Lerninhalte}

\begin{center}
\begin{tikzpicture}[node distance=0.6cm and 1.0cm]
  \node[box] (def) {$Ax = b$: Koeffizientenmatrix,\\erweiterte Matrix $A'$};
  \node[box, below=0.5cm of def] (los) {Lösbarkeit: $\operatorname{Rg}(A) = \operatorname{Rg}(A')$};
  \node[box, below=0.5cm of los] (par) {Partikuläre Lösung $\lambda_0$\\(Merkregel 5.2.9)};
  \node[box, below=0.5cm of par] (hom) {Homogene Lösungsmenge $U$\\(Merkregel 5.2.16)};
  \node[box, below=0.5cm of hom] (str) {$L = \lambda_0 + U$\\Eindeutigkeit $\Leftrightarrow \operatorname{Rg}(A) = n$};
  \draw[arrow] (def) -- (los);
  \draw[arrow] (los) -- (par);
  \draw[arrow] (par) -- (hom);
  \draw[arrow] (hom) -- (str);
\end{tikzpicture}
\end{center}

\subsection*{Lernziele}

Nach Durcharbeiten dieses Kapitels sollten Sie:
\begin{itemize}
  \item entscheiden können, ob ein lineares Gleichungssystem lösbar ist
        (Vergleich der Ränge von $A$ und $A'$),
  \item eine partikuläre Lösung mit der Merkregel~5.2.9 ablesen können,
  \item die Lösungsmenge des zugehörigen homogenen Systems mit Merkregel~5.2.16
        bestimmen können,
  \item die vollständige Lösungsmenge als $\lambda_0 + U$ angeben und wissen,
        wann das System genau eine Lösung hat.
\end{itemize}

\subsection*{Selbstkontrollelement 5}

Lösen Sie das lineare Gleichungssystem über $\mathbb{R}$:
\[
  \begin{pmatrix} 1 & 2 & -1 & 3 \\ 2 & 4 & 0 & 2 \\ 1 & 2 & 1 & -1 \end{pmatrix}
  \begin{pmatrix} x_1 \\ x_2 \\ x_3 \\ x_4 \end{pmatrix}
  = \begin{pmatrix} 1 \\ 2 \\ 0 \end{pmatrix}.
\]

\newpage

% -----------------------------------------------------------------------
\section*{Zielelement 6.1 -- Vektorräume}

\subsection*{Lerninhalte}

\begin{center}
\begin{tikzpicture}[node distance=0.6cm and 1.0cm]
  \node[box] (mot) {Motivation: Vektoren im $\mathbb{R}^2$\\Parallelogramm-Gesetz};
  \node[box, below=0.5cm of mot] (def) {Definition $K$-Vektorraum\\(3 Axiomengruppen)};
  \node[box, below=0.5cm of def] (rec) {Rechenregeln (Prop.\ 6.1.4)\\$0v = 0$,\; $a \cdot 0 = 0$,\; $(-a)v = -av$};
  \node[box, below=0.5cm of rec] (bsp) {Beispiele: $M_{mn}(K)$, $K^n$,\\$K[T]$, Lösungsmenge von $Ax=0$};
  \draw[arrow] (mot) -- (def);
  \draw[arrow] (def) -- (rec);
  \draw[arrow] (rec) -- (bsp);
\end{tikzpicture}
\end{center}

\subsection*{Lernziele}

Nach Durcharbeiten dieses Abschnitts sollten Sie:
\begin{itemize}
  \item die Definition eines $K$-Vektorraums kennen und an neuen Beispielen
        überprüfen können,
  \item die Rechenregeln aus Proposition~6.1.4 aus den Axiomen herleiten können,
  \item die wichtigsten Beispiele für Vektorräume (insbesondere $K^n$, $M_{mn}(K)$
        und Lösungsmengen homogener Gleichungssysteme) kennen.
\end{itemize}

\subsection*{Selbstkontrollelement 6.1}

Ist die Menge $V = \{f : \mathbb{R} \to \mathbb{R} \mid f(0) = 0\}$ mit der üblichen
punktweisen Addition und Skalarmultiplikation ein $\mathbb{R}$-Vektorraum?
Überprüfen Sie die Axiome.

\newpage

% -----------------------------------------------------------------------
\section*{Zielelement 6.2 -- Unterräume}

\subsection*{Lerninhalte}

\begin{center}
\begin{tikzpicture}[node distance=0.6cm and 1.0cm]
  \node[box] (def) {Definition Unterraum\\$U \neq \emptyset$, Axiome vererbt};
  \node[box, below=0.5cm of def] (kri) {Unterraumkriterium (Prop.\ 6.2.3):\\$0 \in U$, $+$-abgeschlossen, $\cdot$-abgeschlossen};
  \node[box, below=0.5cm of kri] (dsc) {Durchschnitt $U \cap W$\\ist Unterraum};
  \node[box, below=0.5cm of dsc] (sum) {Summe $U + W$\\ist Unterraum; $U \cup W$ i.\,A.\ nicht};
  \draw[arrow] (def) -- (kri);
  \draw[arrow] (kri) -- (dsc);
  \draw[arrow] (dsc) -- (sum);
\end{tikzpicture}
\end{center}

\subsection*{Lernziele}

Nach Durcharbeiten dieses Abschnitts sollten Sie:
\begin{itemize}
  \item das Unterraumkriterium anwenden können, um Unterräume zu erkennen oder
        zu widerlegen,
  \item wissen, dass $U \cap W$ und $U + W$ Unterräume sind, $U \cup W$ aber
        im Allgemeinen nicht,
  \item Lösungsmengen homogener linearer Gleichungssysteme als Unterräume von $K^n$
        einordnen können.
\end{itemize}

\subsection*{Selbstkontrollelement 6.2}

Sei $V = \mathbb{R}^3$ und
$U = \left\{ \begin{pmatrix} x \\ y \\ z \end{pmatrix} \in \mathbb{R}^3 \;\middle|\; x + y = 0 \right\}$.
Zeigen Sie mit dem Unterraumkriterium, dass $U$ ein Unterraum von $\mathbb{R}^3$ ist.

\newpage

% -----------------------------------------------------------------------
\section*{Zielelement 6.3 -- Endlich erzeugte Vektorräume}

\subsection*{Lerninhalte}

\begin{center}
\begin{tikzpicture}[node distance=0.6cm and 1.0cm]
  \node[box] (lkb) {Linearkombination\\$a_1 v_1 + \cdots + a_m v_m$};
  \node[box, below=0.5cm of lkb] (erz) {Lineare Hülle $\langle S \rangle$\\ist Unterraum (Prop.\ 6.3.4)};
  \node[box, below=0.5cm of erz] (es)  {Erzeugendensystem:\\$\langle S \rangle = V$};
  \node[box, below=0.5cm of es]  (end) {Endlich erzeugter Vektorraum:\\$V = \langle v_1, \ldots, v_m \rangle$};
  \node[box, below=0.5cm of end] (bsp) {Beispiele/Gegenbeispiel:\\$K^n$, $M_{mn}(K)$, Polynome $\le n$\\vs.\ $K[T]$ nicht endlich erzeugt};
  \draw[arrow] (lkb) -- (erz);
  \draw[arrow] (erz) -- (es);
  \draw[arrow] (es) -- (end);
  \draw[arrow] (end) -- (bsp);
\end{tikzpicture}
\end{center}

\subsection*{Lernziele}

Nach Durcharbeiten dieses Abschnitts sollten Sie:
\begin{itemize}
  \item Linearkombinationen bilden und die lineare Hülle einer Menge beschreiben können,
  \item nachweisen oder widerlegen können, dass eine Menge $S$ ein Erzeugendensystem
        von $V$ ist (durch Lösung eines linearen Gleichungssystems),
  \item endlich erzeugte Vektorräume von nicht endlich erzeugten unterscheiden,
  \item verstehen, warum $K[T]$ nicht endlich erzeugt ist.
\end{itemize}

\subsection*{Selbstkontrollelement 6.3}

Liegt der Vektor $v = \begin{pmatrix} 1 \\ 2 \\ 3 \end{pmatrix}$ in der linearen Hülle von
$v_1 = \begin{pmatrix} 1 \\ 0 \\ 1 \end{pmatrix}$ und $v_2 = \begin{pmatrix} 0 \\ 1 \\ 1 \end{pmatrix}$?
Bestimmen Sie gegebenenfalls die Koeffizienten.

\end{document}
